% Tutorial set: 03
% Source: T3.pdf, Tutorial 3 (Image Enhancement in the Frequency Domain), Questions 1--3
% Reference: T3_Reference solutions.pptx, used only for answer checking
% Session date: 2026-09-08
% Human verified: Yes

\clearpage
\subsection{Tutorial 03：Frequency-domain Image Enhancement}
\label{tutorial:SC4061-TUT03}

\exercisemeta
  {SC4061-TUT03}
  {2026-09-08}
  {T3.pdf，Questions 1--3；T3\_Reference solutions.pptx 仅用于核对}
  {三题均已完成推导、讨论并订正}
  {已确认（2026-09-08）}

\exercisequestion{SC4061-TUT03-Q01}{Spatial Mask、Frequency Response 与 Separability}

\textbf{对应讲义：}
\relatedtopic{SC4061-L04-T03}{FFT 与 Convolution Theorem}；
\relatedtopic{SC4061-L04-T05}{Low-pass、High-pass 与 Notch Filtering}

\subsubsection*{题目}

给定 spatial convolution mask（空间卷积核）
\[
  h(x,y)=\frac15
  \begin{bmatrix}
    0&1&0\\
    1&1&1\\
    0&1&0
  \end{bmatrix}.
\]
（a）利用 translation property（平移性质）推导其 frequency response（频率响应），并解释其 low-pass property（低通性质）。（b）若 $f(x,y)=f_1(x)f_2(y)$，证明二维 Fourier transform 可分离，并说明该性质对 convolution 的意义。

\begin{checkedsolution}
将 mask 写成 shifted impulses（移位冲激）之和：
\[
  h(x,y)=\frac15\bigl[
  \delta(x,y)+\delta(x-1,y)+\delta(x+1,y)
  +\delta(x,y-1)+\delta(x,y+1)\bigr].
\]
由
\[
  f(x-x_0,y-y_0)
  \Longleftrightarrow
  F(u,v)e^{-j2\pi(ux_0/M+vy_0/N)},
\]
得到
\begin{align*}
  H(u,v)
  &=\frac15\left(1+e^{-j2\pi u/M}+e^{j2\pi u/M}
  +e^{-j2\pi v/N}+e^{j2\pi v/N}\right)\\
  &=\boxed{\frac15\left(1+2\cos\frac{2\pi u}{M}
  +2\cos\frac{2\pi v}{N}\right)}.
\end{align*}
特别地，$H(0,0)=1$，所以 DC component（直流分量）完整保留。该 mask 对中心及四邻域求平均；slow variation（缓慢变化）的各项近似同相，而 rapid oscillation（快速振荡）会在求和中相消，因此总体属于 low-pass filter。它并非 ideal、radially monotone low-pass：例如 $H(M/2,N/2)=-3/5$，仍存在 side lobes（旁瓣）与 phase reversal（相位反转）。

对 separable signal（可分离信号），
\begin{align*}
F(u,v)
&=\sum_{y=0}^{N-1}\sum_{x=0}^{M-1}
f_1(x)f_2(y)e^{-j2\pi ux/M}e^{-j2\pi vy/N}\\
&=\left(\sum_{x=0}^{M-1}f_1(x)e^{-j2\pi ux/M}\right)
\left(\sum_{y=0}^{N-1}f_2(y)e^{-j2\pi vy/N}\right)\\
&=\boxed{F_1(u)F_2(v)}.
\end{align*}
若 filter 也满足 $h(x,y)=h_1(x)h_2(y)$，二维 convolution 可依次执行两个 one-dimensional convolutions（一维卷积）。对 $K\times K$ kernel，每个 pixel 的主要乘法量由约 $K^2$ 降为 $2K$。
\end{checkedsolution}

\textbf{检查点：}``averaging mask'' 足以说明总体 smoothing tendency（平滑倾向），但不意味着 $|H(u,v)|$ 沿任意径向都严格单调。

\exercisequestion{SC4061-TUT03-Q02}{Gaussian Fourier Transform 与 Domain Scaling}

\textbf{对应讲义：}
\relatedtopic{SC4061-L04-T02}{DFT、IDFT 与 Spectrum}；
\relatedtopic{SC4061-L04-T05}{Low-pass、High-pass 与 Notch Filtering}

\subsubsection*{题目}

在
\[
  e^{-\pi x^2}\Longleftrightarrow e^{-\pi u^2},
  \qquad
  f(ax)\Longleftrightarrow \frac1{|a|}F\!\left(\frac ua\right)
\]
的 convention 下，证明 standard deviation（标准差）为 $\sigma$ 的 unit-amplitude Gaussian（单位峰值高斯函数）仍变换为 Gaussian，并求其 frequency-domain amplitude 与 standard deviation。

\begin{checkedsolution}
令 $\sigma>0$，将空间域函数配成已知形式：
\[
  e^{-x^2/(2\sigma^2)}=e^{-\pi(ax)^2},
  \qquad
  \pi a^2=\frac1{2\sigma^2},
  \qquad
  a=\frac1{\sqrt{2\pi}\,\sigma}.
\]
使用 scaling theorem：
\begin{align*}
  e^{-x^2/(2\sigma^2)}
  &\Longleftrightarrow
  \frac1a e^{-\pi(u/a)^2}\\
  &=\sqrt{2\pi}\,\sigma\,
  e^{-2\pi^2\sigma^2u^2}\\
  &=\boxed{\sqrt{2\pi}\,\sigma
  \exp\!\left(-\frac{u^2}{2\alpha^2}\right)},
  \qquad
  \boxed{\alpha=\frac1{2\pi\sigma}}.
\end{align*}
前面的 coefficient（系数）给出 frequency-domain amplitude；指数中的 $\alpha$ 给出 width（宽度），两者作用不同。还可用
\[
  F(0)=\int_{-\infty}^{\infty}e^{-x^2/(2\sigma^2)}\,dx
  =\sqrt{2\pi}\,\sigma
\]
检查 amplitude。空间域越宽，频域越窄，并满足
\[
  \boxed{\sigma\alpha=\frac1{2\pi}}.
\]
该常数依赖这里采用的 $e^{-j2\pi ux}$ Fourier convention；更换 convention 会移动 $2\pi$。
\end{checkedsolution}

\textbf{检查点：}本题空间域 Gaussian 没有除以 $\sqrt{2\pi}\sigma$，所以它的面积和 $F(0)$ 都随 $\sigma$ 增大；不能只改指数而漏掉 scaling coefficient。

\exercisequestion{SC4061-TUT03-Q03}{Second Difference 的 Directional High-pass Response}

\textbf{对应讲义：}
\relatedtopic{SC4061-L04-T05}{Low-pass、High-pass 与 Notch Filtering}；
\relatedtopic{SC4061-L05-T04}{Image Edges、Gradient 与 Curvature}

\subsubsection*{题目}

设
\[
  g(x,y)=f(x+1,y)+f(x-1,y)-2f(x,y).
\]
推导 $G(u,v)$ 与 $F(u,v)$ 的关系，判断 filter type（滤波器类型），并说明其方向性。

\begin{checkedsolution}
两次 spatial shift 分别产生相反的 phase factor：
\begin{align*}
  G(u,v)
  &=\left(e^{j2\pi u/M}+e^{-j2\pi u/M}-2\right)F(u,v)\\
  &=\left(2\cos\frac{2\pi u}{M}-2\right)F(u,v)\\
  &=\boxed{-4\sin^2\!\left(\frac{\pi u}{M}\right)F(u,v)}.
\end{align*}
因此 transfer function（传递函数）为
\[
  H(u,v)=-4\sin^2\!\left(\frac{\pi u}{M}\right),
  \qquad
  |H(u,v)|=4\sin^2\!\left(\frac{\pi u}{M}\right).
\]
$H(0,v)=0$，而 $|H(M/2,v)|=4$；它消除沿 $x$ 不变的成分并增强沿 $x$ 的快速变化，所以是 directional high-pass filter（方向性高通滤波器）。负号改变 response sign/phase，不改变 pass type。

原始 DFT array 中 $u=0,1,\ldots,M-1$ 的后半段代表 negative frequencies（负频率）：
\[
  \widetilde u=
  \begin{cases}
    u,&0\le u\le M/2,\\
    u-M,&M/2<u<M.
  \end{cases}
\]
例如 $u=M-1$ 实际是低频 $-1$，而不是更高的正频率。因此响应从 $u=M/2$ 向 $u=M$ 降低，是 frequency wrapping（频率回绕）后的对称部分；在 centred frequency axis 上，$|H|$ 随 $|\widetilde u|$ 增大，在 Nyquist frequency 处最大。

\begin{figure}[htbp]
  \centering
  \resizebox{0.96\textwidth}{!}{\begin{tikzpicture}[x=1.05cm,y=1.0cm,font=\scriptsize]
  \tikzset{
    axis/.style={->,semithick,gray!75!black},
    lowpass/.style={very thick,noteBlue},
    highpass/.style={very thick,ntuRed},
    guide/.style={dashed,gray!55}
  }

  \begin{scope}[shift={(0,0)}]
    \node[font=\small\bfseries] at (3.3,3.65) {Q1: $H(\omega,0)=(3+2\cos\omega)/5$};
    \draw[axis] (0,0) -- (6.65,0) node[right] {$\omega$};
    \draw[axis] (3.3,0) -- (3.3,3.15) node[above] {$H$};
    \draw[guide] (0.25,2.75) -- (6.35,2.75);
    \draw[guide] (0.25,0.99) -- (6.35,0.99);
    \draw[lowpass,smooth,samples=100,domain=-3.14159:3.14159]
      plot ({3.3+0.97*\x},{0.55+2.2*(3+2*cos(\x r))/5});
    \node[left] at (3.3,2.75) {$1$};
    \node[left] at (3.3,0.99) {$1/5$};
    \node[below] at (0.25,0) {$-\pi$};
    \node[below] at (3.3,0) {$0$};
    \node[below] at (6.35,0) {$\pi$};
    \node[align=center,noteBlue] at (3.3,-0.7) {DC passed; high frequencies attenuated};
  \end{scope}

  \begin{scope}[shift={(8.0,0)}]
    \node[font=\small\bfseries] at (3.3,3.65) {Q3: $|H(\omega)|=4\sin^2(\omega/2)$};
    \draw[axis] (0,0) -- (6.65,0) node[right] {$\omega$};
    \draw[axis] (3.3,0) -- (3.3,3.15) node[above] {$|H|$};
    \draw[guide] (0.25,2.75) -- (6.35,2.75);
    \draw[highpass,smooth,samples=100,domain=-3.14159:3.14159]
      plot ({3.3+0.97*\x},{0.55+0.55*4*(sin(\x/2 r))^2});
    \node[left] at (3.3,2.75) {$4$};
    \node[left] at (3.3,0.55) {$0$};
    \node[below] at (0.25,0) {$-\pi$};
    \node[below] at (3.3,0) {$0$};
    \node[below] at (6.35,0) {$\pi$};
    \node[align=center,ntuRed] at (3.3,-0.7) {DC rejected; Nyquist frequencies strongest};
  \end{scope}
\end{tikzpicture}
}
  \caption{Q1 与 Q3 的一维 frequency-response cross-sections（频率响应截面）。横轴使用 centred angular frequency，避免把 raw DFT index 的回绕误认为更高频率。}
  \label{fig:sc4061-tut03-frequency-responses}
\end{figure}

$g$ 是沿 $x$ 的 exact central second difference（精确中心二阶差分）；当 sampling interval 为 $1$ 且 $f$ 足够平滑时，才近似 $\partial^2f/\partial x^2$，不能直接写成严格相等。

方向名称依赖 coordinate convention（坐标约定）：该算子始终检测 intensity 沿 $x$ 变化的边界，edge tangent 与 $x$ 方向垂直。按本课程常用的 image-matrix convention（$x$ 为 row、从上到下），它响应 horizontal edges（水平边缘）；若采用 Cartesian convention（$x$ 从左到右），同一算子则响应 vertical edges（竖直边缘）。
\end{checkedsolution}

\textbf{检查点：}判断 DFT filter type 时看 centred signed frequency（中心化带符号频率）上的 $|H|$，不能把 raw index $u$ 的数值大小直接当作频率高低。

\subsubsection*{本次 Tutorial 收口}

Q1 将 local averaging mask 转为 frequency response，并利用 separability 把二维运算拆成两个一维运算；Q2 展示 Gaussian 在 Fourier transform 下保持 Gaussian 形状，但 amplitude 与 width 按 scaling theorem 同时变化；Q3 把 central second difference 写成 directional high-pass response。复习时优先掌握：shift 对应 phase factor、$G=HF$、空间域与频域 Gaussian 宽度成反比，以及 raw DFT indices 的后半段代表 negative frequencies。
